\chapter{Sleep Paralysis}

\section*{Definition}
Sleep paralysis is a transient, involuntary inability to move voluntary muscles during the transition between wakefulness and sleep (hypnagogic or hypnopompic), lasting seconds to a few minutes. It is typically accompanied by preserved consciousness, awareness of the environment, and often vivid hallucinations.

\section*{Pathophysiology}
Sleep paralysis represents a dissociation of REM sleep features -- specifically the muscle atonia (REM atonia persists into waking) while consciousness and eye movement resume. Brainstem mechanisms (pontine tegmentum, particularly the sublaterodorsal nucleus) inappropriately maintain motor inhibition via spinal motoneuron hyperpolarization.

\section*{Etiology}
\begin{itemize}
  \item \textbf{Primary:} Isolated sleep paralysis (idiopathic), often triggered by sleep disruption
  \item \textbf{Sleep disorders:} Narcolepsy type 1 (with cataplexy) -- sleep paralysis is a core feature, obstructive sleep apnea, insomnia, circadian rhythm disorders
  \item \textbf{Sleep deprivation:} Most common trigger; shift work, jet lag, irregular sleep schedules
  \item \textbf{Psychiatric:} Anxiety disorders, panic disorder, PTSD (especially trauma-related sleep disturbance)
  \item \textbf{Substance-related:} Alcohol use, cannabis use, stimulant withdrawal
  \item \textbf{Genetic:} HLA-DQB1*06:02 allele is strongly associated (especially in narcolepsy)
\end{itemize}

\section*{Indications for Evaluation}
Seek care if:
\begin{itemize}
  \item Episodes are frequent (more than 1--2 per month), distressing, or associated with daytime sleepiness
  \item There are other signs of narcolepsy: cataplexy (sudden muscle weakness triggered by emotion), hypnagogic hallucinations, excessive daytime sleepiness
  \item The individual is fearful of sleeping due to paralysis episodes
\end{itemize}

\section*{Related Symptoms}
\begin{itemize}
  \item Hypnagogic hallucinations, cataplexy, excessive daytime sleepiness, sleep fragmentation, insomnia, narcolepsy, panic during paralysis
\end{itemize}
\textbf{Note:} Up to 40\% of the general population has experienced at least one episode of sleep paralysis; in most cases it is not a sign of a serious disorder.

\section*{Self-Care / Home Management}
\begin{itemize}
  \item Maintain consistent sleep--wake timing and prioritize 7--9 hours of sleep per night
  \item Avoid sleeping on the back (supine position is a risk factor)
  \item Manage stress and practice relaxation techniques before bed
  \item During an episode, try to move small muscles (finger, toe, eye movements) to break the paralysis
\end{itemize}

\section*{Diagnostic Approach}
\begin{itemize}
  \item Clinical history: differentiate from seizures, night terrors, and psychogenic unresponsiveness
  \item Assess for narcolepsy using the Epworth Sleepiness Scale and the narcolepsy symptom questionnaire
  \item Polysomnography with multiple sleep latency test (MSLT) if narcolepsy suspected
  \item HLA typing for DQB1*06:02 if narcolepsy is a consideration
\end{itemize}

\section*{Treatment / Management Overview}
\begin{itemize}
  \item Reassurance and psychoeducation that episodes are benign and self-limited
  \item Optimize sleep hygiene and treat sleep disorders (CPAP for sleep apnea)
  \item For narcolepsy-associated sleep paralysis: antidepressants (SSRIs, SNRIs, TCAs) suppress REM sleep and reduce episodes
  \item Sodium oxybate (consolidates sleep architecture) for narcolepsy with frequent paralysis
\end{itemize}

\clearpage
